% !TEX root = ../main.tex
\begin{ouverture}[avancement]{4}{Avancement et bilan de matière}{Partie 1 · Constitution et transformations de la matière}
{Avancement\enspace·\enspace Tableau d'avancement\enspace·\enspace Réactif limitant\enspace·\enspace État final\enspace·\enspace Transformation totale}
\objectif{Décrire qualitativement l'évolution des quantités de matière au cours d'une transformation.}
\objectif{Établir un tableau d'avancement à partir de l'équation de la réaction et de l'état initial.}
\objectif{Déterminer l'avancement maximal et le réactif limitant ; reconnaître un mélange stœchiométrique.}
\objectif{Déterminer la composition de l'état final et comparer l'avancement final à l'avancement maximal.}
\objectif{Déterminer l'état final d'une transformation totale à l'aide d'un programme Python.}
\prerequis{Relations $n=\dfrac{m}{M}$, $n=c\times V$ et $n=\dfrac{V}{V_\mathrm{m}}$ (chapitre~1).}
\prerequis{Équation de réaction ajustée, réactif limitant (classe de seconde et chapitre~3).}
\end{ouverture}
\essentiel
\essSection{1}{Décrire l'évolution d'un système}
\begin{definition}[Avancement d'une réaction]
Au cours d'une transformation, les réactifs sont consommés et les produits se forment dans les proportions données par les
nombres stœchiométriques de l'équation. L'\textbf{avancement} $x$, en mole, mesure la progression de la réaction : pour la
réaction $a\,\mathrm{A}+b\,\mathrm{B}\longrightarrow c\,\mathrm{C}+d\,\mathrm{D}$,
\[n(\mathrm{A})=n_0(\mathrm{A})-a\,x\qquad n(\mathrm{B})=n_0(\mathrm{B})-b\,x\qquad n(\mathrm{C})=n_0(\mathrm{C})+c\,x\qquad n(\mathrm{D})=n_0(\mathrm{D})+d\,x.\]
\end{definition}
\essSection{2}{Le tableau d'avancement}
\begin{savoirfaire}[Établir un tableau d'avancement]
\begin{questions}
\item Écrire l'équation ajustée de la réaction sur la première ligne.
\item Calculer les quantités de matière initiales et les reporter sur la ligne de l'état initial, où $x=0$.
\item Exprimer, dans l'état intermédiaire, la quantité de chaque espèce en fonction de $x$ : on retranche $a\,x$ pour un
réactif, on ajoute $c\,x$ pour un produit.
\item Compléter l'état final avec la valeur de l'avancement final.
\end{questions}
\end{savoirfaire}
\begin{exemple}
On brûle \qty{0,040}{mol} de méthane dans \qty{0,100}{mol} de dioxygène. La transformation est totale.
\begin{avancement}{4}
\tavEquation{\ce{CH4} & $+\;2\,\ce{O2}$ & $\longrightarrow\;\ce{CO2}$ & $+\;2\,\ce{H2O}$}
\tavTitres{4}
initial & $x=0$ & $0{,}040$ & $0{,}100$ & $0$ & $0$ \\
en cours & $x$ & $0{,}040-x$ & $0{,}100-2x$ & $x$ & $2x$ \\
final & $x_\mathrm{max}=0{,}040$ & $0$ & $0{,}020$ & $0{,}040$ & $0{,}080$ \\
\end{avancement}
\end{exemple}
\schema{04-evolution}
\essSection{3}{Avancement maximal et réactif limitant}
\begin{propriete}[Réactif limitant]
Une transformation totale s'arrête quand un réactif est épuisé : c'est le \textbf{réactif limitant}. L'\textbf{avancement
maximal} $x_\mathrm{max}$ est la plus petite des valeurs de $x$ qui annulent la quantité d'un réactif :
\[x_\mathrm{max}=\text{la plus petite des valeurs }\frac{n_0(\mathrm{A})}{a}\text{ et }\frac{n_0(\mathrm{B})}{b}.\]
Dans l'exemple, $\dfrac{0{,}040}{1}=0{,}040$ et $\dfrac{0{,}100}{2}=0{,}050$ : le méthane est limitant et $x_\mathrm{max}=\qty{0,040}{mol}$.
\end{propriete}
\begin{definition}[Mélange stœchiométrique]
Les réactifs sont introduits dans les \textbf{proportions stœchiométriques} si $\dfrac{n_0(\mathrm{A})}{a}=\dfrac{n_0(\mathrm{B})}{b}$ : ils sont alors tous
épuisés en même temps. Pour brûler \qty{0,040}{mol} de méthane, le mélange stœchiométrique contient \qty{0,080}{mol} de dioxygène.
\end{definition}
\essSection{4}{Transformations totales et non totales}
\begin{definition}[Avancement final]
L'\textbf{avancement final} $x_\mathrm{f}$ se déduit de la composition de l'état final, mesurée expérimentalement. Si $x_\mathrm{f}=x_\mathrm{max}$, la
transformation est \textbf{totale} ; si $x_\mathrm{f}<x_\mathrm{max}$, elle est \textbf{non totale} : tous les réactifs restent présents dans l'état
final. Ainsi, \qty{1,0}{mol} d'acide éthanoïque et \qty{1,0}{mol} d'éthanol ne forment que \qty{0,67}{mol} d'ester : $x_\mathrm{f}=\qty{0,67}{mol}<x_\mathrm{max}=\qty{1,0}{mol}$.
\end{definition}
\begin{python}
# CH4 + 2 O2 -> CO2 + 2 H2O
# chaque espece : [nombre stoechiometrique, quantite initiale en mol]
reactifs = {'CH4': [1, 0.040], 'O2': [2, 0.100]}
produits = {'CO2': [1, 0.0], 'H2O': [2, 0.0]}

candidats = [n0 / a for a, n0 in reactifs.values()]
xmax = min(candidats)
print('xmax =', round(xmax, 4), 'mol')
for espece, (a, n0) in reactifs.items():
    print(espece, round(n0 - a * xmax, 4), 'mol')
for espece, (c, n0) in produits.items():
    print(espece, round(n0 + c * xmax, 4), 'mol')
\end{python}
\begin{console}
xmax = 0.04 mol
CH4 0.0 mol
O2 0.02 mol
CO2 0.04 mol
H2O 0.08 mol
\end{console}
\noindent\capacite{Capacité numérique} Pour chaque réactif, la liste \code{candidats} contient la valeur de $x$ qui l'épuiserait ; la plus petite
est $x_\mathrm{max}$. Les deux boucles calculent ensuite les quantités de matière de l'état final, $n_0-a\,x_\mathrm{max}$ pour un réactif et
$n_0+c\,x_\mathrm{max}$ pour un produit. Il suffit de modifier les deux dictionnaires pour traiter une autre transformation totale.
\begin{attention}
Pour trouver le réactif limitant, on ne compare pas directement les quantités initiales : on compare les quotients
$\dfrac{n_0}{a}$, qui tiennent compte des nombres stœchiométriques. Dans l'exemple, il y a plus de dioxygène que de méthane, mais
chaque molécule de méthane en consomme deux.
\end{attention}
\begin{experience}[Des ballons qui gonflent]
Quatre flacons contiennent chacun \qty{50}{mL} d'acide chlorhydrique à \qty{0,20}{mol.L^{-1}}. On y fait tomber \qty{0,05}{g}, \qty{0,10}{g}, \qty{0,15}{g}
et \qty{0,20}{g} de magnésium, puis on coiffe chaque flacon d'un ballon. Le dihydrogène formé gonfle les ballons, de plus en plus
jusqu'au troisième ; les deux derniers ont le même volume, environ \qty{120}{mL} : à partir de là, l'acide est le réactif limitant.
\end{experience}
\begin{rediger}
Écrire l'équation, calculer les quantités initiales avec leurs unités, dresser le tableau d'avancement, comparer les
quotients $\dfrac{n_0}{a}$ pour conclure sur le réactif limitant, puis donner la composition de l'état final en quantités de
matière, et en masses ou en volumes si l'énoncé le demande.
\end{rediger}
\exercices

\rubrique{Découverte}
\begin{blocLivre}{ex-4-1}
\begin{exercice}[1]{Compter des molécules}
Le schéma représente un mélange de dihydrogène et de dioxygène avant et après la réaction $\ce{2H2 + O2 -> 2H2O}$.
\schema{04-particules}
\begin{questions}
\item Combien de molécules de chaque réactif l'état initial contient-il ?
\item Combien de fois la réaction peut-elle se produire au maximum ? Quel réactif est entièrement consommé ?
\item Vérifier que l'état final représenté conserve les atomes de chaque élément.
\item Un échantillon réel contient environ $10^{23}$ molécules. Pourquoi exprime-t-on l'avancement en moles ?
\end{questions}
\end{exercice}
\end{blocLivre}
\begin{blocLivre}{ex-4-2}
\begin{exercice}[1]{Le fer et le soufre}
On chauffe un mélange de \qty{0,100}{mol} de fer et de \qty{0,080}{mol} de soufre. Il se forme du sulfure de fer selon la réaction
totale d'équation $\ce{Fe + S -> FeS}$.
\begin{donnees}
Masses molaires (\unit{g.mol^{-1}}) : Fe 55,8 ; S 32,1 ; FeS 87,9.
\end{donnees}
\begin{questions}
\item Calculer les masses de fer et de soufre du mélange initial.
\item Établir le tableau d'avancement de la transformation.
\item Déterminer l'avancement maximal et le réactif limitant.
\item Calculer les masses des espèces présentes dans l'état final et vérifier la conservation de la masse.
\end{questions}
\end{exercice}
\end{blocLivre}
\begin{blocLivre}{ex-4-3}
\begin{exercice}[1]{Réactif limitant ou mélange stœchiométrique ?}
Un ruban de magnésium brûle dans le dioxygène avec une flamme éblouissante, selon la réaction totale d'équation
$\ce{2Mg + O2 -> 2MgO}$. On étudie trois mélanges.
\par\smallskip\noindent\hfill\begin{tabular}{lccc}
\toprule
Mélange & A & B & C \\
\midrule
$n_0(\ce{Mg})$ (\unit{mol}) & 0,20 & 0,20 & 0,30 \\
$n_0(\ce{O2})$ (\unit{mol}) & 0,10 & 0,20 & 0,10 \\
\bottomrule
\end{tabular}\hfill\null\par\smallskip
\begin{questions}
\item Pour chaque mélange, calculer les quotients $\dfrac{n_0(\ce{Mg})}{2}$ et $\dfrac{n_0(\ce{O2})}{1}$.
\item En déduire l'avancement maximal et le réactif limitant de chaque mélange.
\item Quel mélange est stœchiométrique ? Écrire la relation que vérifient alors les quantités initiales.
\item Donner la composition de l'état final du mélange B.
\end{questions}
\end{exercice}
\end{blocLivre}
\begin{blocLivre}{ex-4-4}
\begin{exercice}[1]{Brûler du propane}
Un réchaud brûle \qty{8,8}{g} de propane $\ce{C3H8}$. La combustion complète a pour équation $\ce{C3H8 + 5O2 -> 3CO2 + 4H2O}$.
\begin{donnees}
$M(\ce{C3H8})=\qty{44,0}{g.mol^{-1}}$ ; $M(\ce{CO2})=\qty{44,0}{g.mol^{-1}}$ ; $V_\mathrm{m}=\qty{24,0}{L.mol^{-1}}$ à \qty{20}{\degreeCelsius}. L'air contient 21~\% de dioxygène en volume.
\end{donnees}
\begin{questions}
\item Calculer la quantité de matière de propane.
\item Quelle quantité de dioxygène faut-il pour un mélange stœchiométrique ? Quel volume occupe-t-elle à \qty{20}{\degreeCelsius} ?
\item En déduire le volume d'air nécessaire.
\item Calculer la masse de dioxyde de carbone formée.
\end{questions}
\end{exercice}
\end{blocLivre}
\begin{blocLivre}{ex-4-5}
\begin{exercice}[1]{Une transformation non totale}
On mélange \qty{1,0}{mol} d'acide éthanoïque et \qty{1,0}{mol} d'éthanol. Il se forme de l'éthanoate d'éthyle, un solvant à l'odeur de
dissolvant pour vernis, et de l'eau : $\ce{CH3COOH + C2H6O -> C4H8O2 + H2O}$. Après quelques jours, la composition ne
change plus ; le système contient alors \qty{0,67}{mol} d'ester.
\begin{questions}
\item Établir le tableau d'avancement et déterminer l'avancement maximal.
\item Déterminer l'avancement final.
\item La transformation est-elle totale ? Justifier.
\item Donner la composition de l'état final, puis la masse d'ester obtenue ($M=\qty{88,0}{g.mol^{-1}}$).
\end{questions}
\end{exercice}
\end{blocLivre}
\rubrique{Entraînement}
\begin{blocLivre}{ex-4-6}
\begin{exercice}[2]{Un programme pour l'état final}
On reprend le programme de l'essentiel, qui détermine l'état final de la combustion du méthane.
\begin{questions}
\item \capacite{Capacité numérique}Que contient la liste \code{candidats} dans l'exemple ? Pourquoi \code{xmax} est-il le minimum de cette liste ?
\item On remplace la quantité initiale de dioxygène par \num{0.060}. Prévoir les cinq lignes affichées.
\item Écrire les dictionnaires à utiliser pour la réaction $\ce{Fe + S -> FeS}$, avec \qty{0,100}{mol} de fer et \qty{0,080}{mol} de soufre.
Prévoir l'affichage.
\item Proposer une ligne à ajouter dans la première boucle pour afficher le nom du réactif limitant.
\end{questions}
\end{exercice}
\end{blocLivre}
\begin{blocLivre}{ex-4-7}
\begin{exercice}[2]{Le coussin gonflable de sécurité}
En cas de choc, un capteur déclenche dans le coussin gonflable (airbag) la décomposition de l'azoture de sodium $\ce{NaN3}$ :
$\ce{2NaN3 -> 2Na + 3N2}$. Le diazote formé gonfle le coussin en quelques centièmes de seconde.
\begin{donnees}
Masses molaires (\unit{g.mol^{-1}}) : N 14,0 ; O 16,0 ; Na 23,0 ; K 39,1. $V_\mathrm{m}=\qty{24,0}{L.mol^{-1}}$.
\end{donnees}
\begin{questions}
\item Le coussin contient \qty{130}{g} d'azoture de sodium. Calculer sa quantité de matière et l'avancement maximal.
\item Calculer le volume de diazote libéré à \qty{20}{\degreeCelsius}. Comparer au volume d'un coussin conducteur, d'environ \qty{60}{L}.
\item Quelle masse de sodium se forme-t-elle ?
\item Le sodium, dangereux, est éliminé par du nitrate de potassium : $\ce{10Na + 2KNO3 -> K2O + 5Na2O + N2}$. Quelle masse de nitrate
de potassium faut-il prévoir ?
\end{questions}
\end{exercice}
\end{blocLivre}
\begin{blocLivre}{ex-4-8}
\begin{exercice}[2]{Souder des rails}
Pour souder deux rails, on enflamme dans un moule un mélange d'aluminium en poudre et d'oxyde de fer(III) : la réaction,
très exothermique, produit du fer fondu qui remplit l'espace entre les rails. Son équation est $\ce{2Al + Fe2O3 -> Al2O3 + 2Fe}$.
\begin{donnees}
Masses molaires (\unit{g.mol^{-1}}) : O 16,0 ; Al 27,0 ; Fe 55,8.
\end{donnees}
\begin{questions}
\item Une charge contient \qty{2,70}{kg} d'aluminium et \qty{7,98}{kg} d'oxyde de fer(III). Calculer les quantités de matière.
\item Montrer que le mélange est stœchiométrique.
\item Calculer la masse de fer produite.
\item Calculer la masse d'oxyde d'aluminium formée et vérifier la conservation de la masse.
\end{questions}
\end{exercice}
\end{blocLivre}
\begin{blocLivre}{ex-4-9}
\begin{exercice}[2]{Vinaigre et bicarbonate}
On verse \qty{2,10}{g} d'hydrogénocarbonate de sodium, $\ce{NaHCO3}$, dans \qty{50,0}{mL} de vinaigre contenant \qty{50}{g} d'acide éthanoïque
par litre. La réaction a pour équation $\ce{CH3COOH + HCO3- -> CH3COO- + CO2 + H2O}$.
\schema{04-ex09}
\begin{donnees}
Masses molaires (\unit{g.mol^{-1}}) : $\ce{NaHCO3}$ 84,0 ; $\ce{CH3COOH}$ 60,0 ; $\ce{CO2}$ 44,0.
\end{donnees}
\begin{questions}
\item Calculer les quantités initiales des deux réactifs.
\item Déterminer l'avancement maximal et le réactif limitant. Quelle masse de dioxyde de carbone peut-on attendre ?
\item Déduire des mesures l'avancement final.
\item Comparer $x_\mathrm{f}$ et $x_\mathrm{max}$. Comment expliquer le petit écart, sachant que le dioxyde de carbone est un peu soluble dans l'eau ?
\end{questions}
\end{exercice}
\end{blocLivre}
\begin{blocLivre}{ex-4-10}
\begin{exercice}[2]{Un précipité blanc}
On mélange \qty{20,0}{mL} de nitrate d'argent, $\ce{Ag+}+\ce{NO3-}$, à \qty{0,050}{mol.L^{-1}} et \qty{10,0}{mL} de chlorure de sodium,
$\ce{Na+}+\ce{Cl-}$, à \qty{0,080}{mol.L^{-1}}. Il se forme un précipité blanc de chlorure d'argent : $\ce{Ag+ + Cl- -> AgCl}$, réaction totale.
\begin{donnees}
$M(\ce{AgCl})=\qty{143,4}{g.mol^{-1}}$.
\end{donnees}
\begin{questions}
\item Calculer les quantités initiales d'ions argent et d'ions chlorure.
\item Établir le tableau d'avancement et identifier le réactif limitant.
\item Calculer la masse de précipité formé.
\item Calculer les concentrations des ions encore présents dans la solution finale, de volume \qty{30,0}{mL}.
\end{questions}
\end{exercice}
\end{blocLivre}
\begin{blocLivre}{ex-4-11}
\begin{exercice}[2]{Retrouver une quantité initiale}
On ajoute une masse inconnue $m$ de zinc en poudre à \qty{50,0}{mL} d'une solution de sulfate de cuivre à \qty{0,100}{mol.L^{-1}}. Une
fois le zinc entièrement consommé, une mesure d'absorbance donne $[\ce{Cu^{2+}}]=\qty{0,040}{mol.L^{-1}}$. L'équation de la réaction est
$\ce{Zn + Cu^{2+} -> Zn^{2+} + Cu}$.
\begin{questions}
\item Établir le tableau d'avancement en notant $n_0$ la quantité initiale de zinc.
\item Exprimer la quantité finale d'ions cuivre(II) en fonction de $x_\mathrm{f}$, puis résoudre cette équation du premier degré.
\item En déduire $n_0$ et la masse $m$ de zinc ($M=\qty{65,4}{g.mol^{-1}}$).
\item Calculer la concentration finale en ions zinc(II).
\end{questions}
\end{exercice}
\end{blocLivre}
\rubrique{Objectif DS}
\begin{blocLivre}{ex-4-12}
\begin{exercice}[3]{Un réchaud de camping}
\begin{docu}{1}{La cartouche}
Une cartouche de réchaud contient \qty{190}{g} de butane $\ce{C4H10}$. Sa combustion complète a pour équation
$\ce{2C4H10 + 13O2 -> 8CO2 + 10H2O}$.
\end{docu}
\begin{docu}{2}{Une mise en garde}
Le fabricant précise : « Ne pas utiliser dans un espace clos, une tente ou un véhicule. » Lorsque le dioxygène vient à
manquer, la combustion devient incomplète et produit du monoxyde de carbone, un gaz incolore, inodore et mortel.
\end{docu}
\begin{donnees}
Masses molaires (\unit{g.mol^{-1}}) : H 1,0 ; C 12,0 ; O 16,0. $V_\mathrm{m}=\qty{24,0}{L.mol^{-1}}$ ; l'air contient 21~\% de dioxygène en volume.
\end{donnees}
\begin{questions}
\item Calculer la quantité de butane de la cartouche et l'avancement maximal de sa combustion complète.
\item Calculer le volume de dioxygène, puis le volume d'air, nécessaires pour brûler toute la cartouche.
\item Calculer les masses de dioxyde de carbone et d'eau produites.
\item \initiative Une tente de \qty{3}{m^3} contient environ \qty{0,6}{m^3} de dioxygène. Expliquer, à l'aide d'un raisonnement sur le réactif
limitant, la mise en garde du document~2.
\end{questions}
\end{exercice}
\end{blocLivre}
\begin{blocLivre}{ex-4-13}
\begin{exercice}[3]{Fabriquer de l'ammoniac}
\begin{docu}{1}{Le procédé Haber-Bosch}
L'ammoniac $\ce{NH3}$, matière première des engrais azotés, est synthétisé à partir de diazote et de dihydrogène :
$\ce{N2 + 3H2 -> 2NH3}$. Dans le réacteur, à \qty{450}{\degreeCelsius} et sous \qty{200}{bar}, la transformation n'est pas totale.
\end{docu}
\begin{docu}{2}{Un essai}
On introduit \qty{1,00}{mol} de diazote et \qty{3,00}{mol} de dihydrogène. À la sortie du réacteur, le mélange contient \qty{0,30}{mol}
d'ammoniac. On le refroidit : l'ammoniac se liquéfie et est séparé, tandis que les gaz restants sont renvoyés dans le
réacteur.
\end{docu}
\begin{questions}
\item Le mélange initial est-il stœchiométrique ? Déterminer l'avancement maximal.
\item Établir le tableau d'avancement et déterminer l'avancement final.
\item La transformation est-elle totale ? Donner la composition du mélange sortant du réacteur.
\item \initiative À chaque passage, \qty{15}{\percent} du diazote présent est transformé. Combien de passages faut-il pour transformer au
moins \qty{90}{\percent} du diazote initial ?
\end{questions}
\end{exercice}
\end{blocLivre}
\begin{blocLivre}{ex-4-14}
\begin{exercice}[3]{La lampe à carbure}
\begin{docu}{1}{L'éclairage des spéléologues}
Une lampe à carbure contient deux réservoirs. L'eau du réservoir supérieur goutte sur du carbure de calcium $\ce{CaC2}$ : il
se forme de l'acétylène $\ce{C2H2}$, un gaz qui brûle à la sortie d'un bec, et de l'hydroxyde de calcium :
\[\ce{CaC2 + 2H2O -> C2H2 + Ca(OH)2}\qquad\qquad\ce{2C2H2 + 5O2 -> 4CO2 + 2H2O}.\]
\end{docu}
\begin{docu}{2}{Une charge}
Une charge contient \qty{100}{g} de carbure de calcium. Le bec débite environ \qty{12}{L} d'acétylène par heure.
\end{docu}
\begin{donnees}
Masses molaires (\unit{g.mol^{-1}}) : H 1,0 ; C 12,0 ; O 16,0 ; Ca 40,1. $V_\mathrm{m}=\qty{24,0}{L.mol^{-1}}$ (température de la grotte assimilée à \qty{20}{\degreeCelsius}).
\end{donnees}
\begin{questions}
\item Calculer la quantité de carbure de calcium d'une charge.
\item Quelle masse minimale d'eau faut-il prévoir pour la transformer entièrement ? Quel volume d'acétylène obtient-on ?
\item Quel volume de dioxygène la combustion de cet acétylène consomme-t-elle ?
\item \initiative Estimer l'autonomie de la lampe avec une charge.
\end{questions}
\end{exercice}
\end{blocLivre}
\begin{blocLivre}{ex-4-15}
\begin{exercice}[3]{La paille de fer qui brûle}
\begin{docu}{1}{L'expérience}
De la paille de fer chauffée au rouge est plongée dans un flacon de dioxygène : elle brûle en projetant des étincelles et
se transforme en oxyde magnétique $\ce{Fe3O4}$. La réaction est totale.
\end{docu}
\begin{docu}{2}{Un graphique}
Les trois droites représentent les quantités de matière du fer, du dioxygène et de l'oxyde en fonction de l'avancement.
\schema{04-ex15}
\end{docu}
\begin{donnees}
Masses molaires (\unit{g.mol^{-1}}) : O 16,0 ; Fe 55,8. $V_\mathrm{m}=\qty{24,0}{L.mol^{-1}}$.
\end{donnees}
\begin{questions}
\item Quelle droite correspond au produit ? Justifier.
\item \initiative À l'aide des pentes des droites, identifier les deux autres espèces et établir l'équation de la réaction.
\item Lire les quantités initiales ; en déduire la masse de fer et le volume de dioxygène introduits.
\item Déterminer l'avancement maximal, le réactif limitant et la masse d'oxyde magnétique obtenue.
\end{questions}
\end{exercice}
\end{blocLivre}
